\phantomsection
\section*{Глава 3. Разработка веб-приложения}
\addcontentsline{toc}{section}{Глава 3. Разработка веб-приложения}
\stepcounter{section}

В третьей главе описывается программная реализация веб-приложения. Кодовая база организована как Django-проект \texttt{pcshop}, разделённый на четыре приложения (Django apps): \texttt{accounts} (пользователи, аутентификация и аналитика), \texttt{catalog} (категории и комплектующие), \texttt{builds} (сборки ПК, состав и конфигуратор), \texttt{orders} (корзина, заказы и статусы). Реализация выполнялась на основе официальной документации Django~\cite{django_docs} и учебных материалов~\cite{hawkins2022django}.

\subsection{Реализация базовой структуры приложения и шаблонов страниц}

Базовая структура проекта построена в соответствии с принципом «приложение на смысловую область»: каждая Django-app содержит модели, миграции, представления, формы, маршруты и админ-регистрацию своей сущности. Глобальные настройки (\texttt{settings.py}), маршрутизация (\texttt{urls.py}) и WSGI-точки входа собраны в каталоге \texttt{pcshop/}. В качестве СУБД по умолчанию используется SQLite (без дополнительной настройки), переключение на PostgreSQL выполняется через переменные окружения в файле \texttt{.env} с помощью пакета \texttt{python-decouple}.

Структура каталогов проекта приведена в таблице~\ref{tab:project-layout}.

\begin{table}[ht!]
\caption{Структура каталогов проекта \texttt{pcshop}}
\label{tab:project-layout}
\small
\begin{tabularx}{\textwidth}{|>{\RaggedRight\arraybackslash}p{0.30\textwidth}|>{\RaggedRight\arraybackslash}X|}
\hline
\textbf{Каталог / файл} & \textbf{Назначение} \\ \hline
\texttt{pcshop/} & Корневой пакет проекта: \texttt{settings.py}, \texttt{urls.py}, \texttt{views.py}, \texttt{wsgi.py} \\ \hline
\texttt{accounts/} & Пользователь, авторизация, ролевые миксины, аналитика, дашборд \\ \hline
\texttt{catalog/} & Категории, комплектующие, импорт/экспорт CSV \\ \hline
\texttt{builds/} & Сборки ПК, состав, конфигуратор \\ \hline
\texttt{orders/} & Корзина, заказы, позиции, статусы, контекст-процессор корзины \\ \hline
\texttt{templates/} & HTML-шаблоны Bootstrap~5, сгруппированные по приложениям \\ \hline
\texttt{fixtures/} & Стартовые данные (статусы заказов, категории, образцы товаров) \\ \hline
\end{tabularx}
\end{table}

\textbf{Клиентская часть} реализована на CSS-фреймворке Bootstrap~5 (подключается через CDN). Разработан единый базовый шаблон \texttt{templates/base.html}, обеспечивающий:
\begin{itemize}
    \item стек технологий Google Fonts (Inter для текста, JetBrains Mono для кода);
    \item «стеклянную» липкую навигационную панель с эффектом \texttt{backdrop-filter: blur} и адаптивным меню по ролям пользователя;
    \item единую палитру (основной цвет --- индиго \texttt{\#4F46E5}, акцент --- янтарный \texttt{\#F59E0B}), скругления, мягкие тени, hover-эффекты;
    \item современную типографику и адаптивную сетку для мобильных устройств.
\end{itemize}

Главная страница \texttt{templates/home.html} содержит hero-блок с градиентным фоном и призывом к действию, сетку категорий с иконками \texttt{Bootstrap Icons}, блок популярных сборок и блок преимуществ. Все вторичные шаблоны (каталог, корзина, личный кабинет, аналитика и др.) наследуются от базового через \texttt{\{\% extends \%\}}.

Для отображения числа позиций корзины в навигационной панели реализован контекст-процессор \texttt{orders/context\_processors.py}, зарегистрированный в \texttt{TEMPLATES[...]['OPTIONS']['context\_processors']}. Процессор читает данные только из сессии браузера (без обращений к базе данных) и добавляет переменную \texttt{cart\_count} в контекст каждого шаблона. На иконке «Корзина» отображается оранжевый числовой бейдж при наличии позиций.

\subsection{Реализация аутентификации и проверки прав доступа}

Для разграничения доступа реализована собственная модель пользователя \texttt{accounts.User}, наследующаяся от \texttt{AbstractUser}. Логин по умолчанию заменён на email: установлены \texttt{USERNAME\_FIELD = 'email'} и \texttt{REQUIRED\_FIELDS = []}, поле \texttt{username} отключено. Модель содержит поле \texttt{role} (\texttt{TextChoices}: \texttt{client}, \texttt{manager}, \texttt{admin}), а также поля \texttt{full\_name} и \texttt{phone}. Реализован собственный \texttt{UserManager}, корректно создающий обычных пользователей и суперпользователей с автоматическим назначением роли \texttt{admin}.

\textbf{Формы аутентификации.} Реализованы три ключевые формы:
\begin{itemize}
    \item \texttt{RegisterForm} (наследник \texttt{UserCreationForm}) --- регистрация клиента; роль принудительно выставляется в \texttt{client}; CSS-классы Bootstrap присваиваются автоматически в \texttt{\_\_init\_\_};
    \item \texttt{EmailAuthenticationForm} (наследник \texttt{AuthenticationForm}) --- замена \texttt{username} на \texttt{EmailField};
    \item \texttt{ProfileForm} и \texttt{BootstrapPasswordChangeForm} --- редактирование профиля и смена пароля.
\end{itemize}

\textbf{Представления.} Аутентификация реализована через class-based views: \texttt{RegisterView} (на базе \texttt{CreateView} с автоматическим логином после регистрации), \texttt{EmailLoginView}, \texttt{CustomLogoutView}, \texttt{ProfileView}, \texttt{CustomPasswordChangeView}.

\textbf{Ролевая модель доступа.} Модуль \texttt{accounts/permissions.py} содержит базовый миксин \texttt{RoleRequiredMixin}, наследующийся от \texttt{LoginRequiredMixin} и \texttt{UserPassesTestMixin}. Базовый миксин проверяет роль пользователя и предоставляет доступ суперпользователям и пользователям с ролью \texttt{admin} безусловно. От него наследуются специализированные миксины \texttt{ClientRequiredMixin}, \texttt{ManagerRequiredMixin}, \texttt{AdminRequiredMixin} и \texttt{StaffRequiredMixin} (менеджер или администратор). Для функциональных view реализованы аналогичные декораторы (\texttt{@manager\_required}, \texttt{@staff\_required} и др.). Реализация полностью соответствует матрице прав доступа, спроектированной в подразделе~2.1.

\subsection{Реализация CRUD-интерфейсов}

CRUD-операции над основными сущностями реализованы в двух плоскостях: административная панель Django для технического сопровождения и публичные/менеджерские страницы на Bootstrap~5 для основной работы пользователей.

\textbf{Каталог комплектующих.} В \texttt{catalog/views.py} реализованы:
\begin{itemize}
    \item \texttt{CatalogListView} --- публичная страница каталога с фильтрацией (поиск по названию и бренду, выбор категории, диапазон цен, наличие на складе) и пагинацией по 12 товаров;
    \item \texttt{ComponentDetailView} --- карточка товара с характеристиками из JSON-поля \texttt{specs}, ценой, остатком и кнопкой «В корзину»;
    \item \texttt{ComponentManageListView} (с миксином \texttt{StaffRequiredMixin}) --- менеджерская таблица товаров с фильтрами «низкий остаток» и категориями;
    \item \texttt{ComponentCreateView}, \texttt{ComponentUpdateView}, \texttt{ComponentDeleteView} --- стандартные CBV для CRUD с формой \texttt{ComponentForm} на Bootstrap.
\end{itemize}

\textbf{Сборки ПК.} В \texttt{builds/views.py} реализованы публичный список (\texttt{BuildListView}), детальная страница со составом сборки (\texttt{BuildDetailView}) и менеджерский CRUD с inline-formset. Класс \texttt{BuildEditView} использует \texttt{BuildItemFormSet} (\texttt{inlineformset\_factory}), что позволяет редактировать состав сборки на одной странице с самой сборкой; сохранение выполняется атомарно в одной транзакции.

\textbf{Заказы.} В \texttt{orders/views.py} реализованы:
\begin{itemize}
    \item \texttt{OrderListView} --- история заказов клиента с пагинацией;
    \item \texttt{OrderDetailView} --- детальная страница заказа с проверкой прав (клиент видит только свои заказы, менеджер и админ --- любые);
    \item \texttt{OrderManageListView} --- менеджерская таблица заказов с поиском и фильтрами по статусу;
    \item функциональное представление \texttt{manage\_change\_status} --- смена статуса заказа менеджером.
\end{itemize}

Особого внимания заслуживает функция \texttt{\_create\_order}: списание остатков реализовано через выражение \texttt{F('stock')}, что гарантирует атомарность операции на уровне СУБД и полностью исключает состояние гонки (race condition) при одновременном оформлении нескольких заказов одного товара:

\begin{lstlisting}[language=Python, caption={Атомарное списание остатков через F()-выражение}]
Component.objects.filter(pk=pk).update(stock=F('stock') - need)
\end{lstlisting}

Всё создание заказа обёрнуто в \texttt{transaction.atomic()}: при любой ошибке (нехватка остатков, ошибка записи) все изменения откатываются целиком.

\textbf{Личный кабинет (дашборд).} Представление \texttt{DashboardView} динамически формирует контекст в зависимости от роли пользователя: для клиента передаются последние 5 заказов с \texttt{select\_related('status')} (без N+1-запросов к статусам); для менеджера и администратора --- последние 10 заказов всех клиентов. Шаблон \texttt{accounts/dashboard.html} отображает: для клиента --- KPI-карточку с числом заказов, ссылки на конфигуратор и каталог, таблицу истории заказов; для менеджера --- карточки быстрого доступа к каталогу, сборкам, заказам и аналитике; для администратора --- дополнительно ссылки на управление пользователями и экспорт/импорт CSV.

\textbf{Каталог --- улучшения UX.} Страница каталога дополнена счётчиком найденных позиций в заголовке, боковой панелью быстрой фильтрации по категориям с подсветкой активной и кнопкой «Сбросить» фильтры. Пагинация реализована с сохранением GET-параметров --- переход на следующую страницу не сбрасывает выбранные фильтры (все параметры передаются через цикл в URL пагинационных ссылок).

\subsection{Реализация модуля сборки ПК из комплектующих (конфигуратор)}

Модуль конфигуратора (URL \texttt{/builds/configure/}) реализован как однострничное приложение на Alpine.js поверх Django-шаблона, что позволяет получить полностью реактивный пользовательский интерфейс без отдельной сборки фронтенда. Архитектурное решение --- frontend получает все доступные комплектующие в момент загрузки страницы (через \texttt{json\_script}) и далее работает без обращений к серверу, кроме финального сохранения сборки.

\textbf{Архитектура шагов.} В \texttt{builds/views.py} определена константа \texttt{CONFIGURATOR\_STEPS} --- список восьми шагов: процессор, материнская плата, оперативная память, накопитель, видеокарта, блок питания, корпус, охлаждение. Каждый шаг имеет поля \texttt{key}, \texttt{slug} категории, \texttt{title}, имя иконки \texttt{Bootstrap Icons} и флаг \texttt{required}. В контекст шаблона передаётся уже сгруппированный по шагам JSON со списком всех активных комплектующих с положительным остатком.

\textbf{Серверная часть} конфигуратора состоит из двух представлений:
\begin{itemize}
    \item \texttt{ConfiguratorView} (\texttt{TemplateView}) --- отдаёт страницу с шагами и JSON-данными комплектующих, отсортированных по цене;
    \item функциональное представление \texttt{configurator\_save} (\texttt{POST}, \texttt{@login\_required}) --- принимает JSON со словарём \texttt{selection} (ключ шага $\rightarrow$ id комплектующего), валидирует, создаёт \texttt{Build} и \texttt{BuildItem}'ы атомарно (\texttt{transaction.atomic}) и опционально добавляет сборку в корзину пользователя.
\end{itemize}

\textbf{Клиентская часть} реализована на Alpine.js (подключён через CDN, ${\sim}$15~КБ) с описанной далее логикой.

\emph{Реактивные компоненты:}
\begin{itemize}
    \item \texttt{steps} --- массив шагов с комплектующими, инициализируется из \texttt{json\_script};
    \item \texttt{currentStep} --- индекс текущего шага мастера;
    \item \texttt{selection} --- словарь выбранных комплектующих (синхронизируется с \texttt{localStorage});
    \item вычисляемые методы \texttt{filteredComponents()}, \texttt{isCompatible()}, \texttt{estimatedPower()}, \texttt{totalPrice()}, \texttt{canCheckout()}.
\end{itemize}

\emph{Правила совместимости} реализованы в методе \texttt{isCompatible(component)} и проверяются динамически при отрисовке каждой карточки:
\begin{enumerate}
    \item \textbf{Процессор $\leftrightarrow$ материнская плата}: сокеты должны совпадать (\texttt{cpu.specs.sockets == mb.specs.socket}).
    \item \textbf{Материнская плата $\leftrightarrow$ память}: тип памяти должен совпадать (\texttt{mb.specs.memory == ram.specs.type}; например, DDR5).
    \item \textbf{Блок питания}: расчётное энергопотребление (сумма TDP процессора, TDP видеокарты и константы 80 Вт на остальные узлы) не должно превышать мощность БП. При превышении карточка БП помечается как несовместимая, при разности менее 100 Вт выводится предупреждение «БП на пределе».
    \item \textbf{Кулер $\leftrightarrow$ процессор}: расчётный TDP кулера должен покрывать TDP процессора.
\end{enumerate}

Несовместимые позиции автоматически помечаются красным бейджем «несовместимо», приглушаются (\texttt{opacity: 0.5}) и блокируются от клика. На самой дешёвой совместимой позиции каждого шага автоматически отображается бейдж «рекомендуем».

\emph{Визуальные элементы интерфейса:}
\begin{itemize}
    \item \textbf{Пошаговый стайплер} с горизонтальной прокруткой и цветовой индикацией состояния (пустой, активный, выполненный);
    \item \textbf{Градиентный прогресс-бар} (индиго $\rightarrow$ фиолетовый $\rightarrow$ розовый), показывающий долю выполненных шагов;
    \item \textbf{Реактивная сводка справа} (\texttt{position: sticky}) с иконками шагов, выбранными комплектующими и ценами;
    \item \textbf{Индикатор мощности БП} с цветной полосой (зелёный $\rightarrow$ жёлтый $\rightarrow$ красный) и текстовой подсказкой о требуемой мощности;
    \item \textbf{Алерты совместимости} в верхней части страницы при конфликтах между выбранными элементами;
    \item \textbf{Сохранение в \texttt{localStorage}} --- при возврате на страницу пользователь видит свой предыдущий выбор.
\end{itemize}

\emph{Завершение сборки.} После выбора всех обязательных шагов разблокируются кнопки «В корзину» (с указанной суммой) и «Сохранить сборку». При нажатии выполняется \texttt{fetch}-запрос на \texttt{/builds/configure/save/} с CSRF-токеном из cookie; сервер создаёт сборку и возвращает URL для редиректа (страница сборки или корзина).

\subsection{Реализация агрегирования данных по остаткам, продажам и прибыли}

Аналитический модуль реализован в виде представления \texttt{AnalyticsView} (\texttt{accounts/views.py}) с использованием миксина \texttt{StaffRequiredMixin} (доступ только менеджерам и администраторам). Все расчёты выполняются на стороне СУБД через агрегирующие методы Django ORM, что обеспечивает приемлемую производительность даже при больших объёмах заказов.

\textbf{Период выборки} задаётся параметром \texttt{?days=N} (по умолчанию 30 дней; поддерживаются значения 7, 30, 90 и 365). Все выборки фильтруются по \texttt{created\_at}, отменённые заказы (статус \texttt{cancelled}) исключаются.

\textbf{Ключевые показатели} рассчитываются за один запрос с использованием выражений \texttt{Sum} и \texttt{F}:
\begin{itemize}
    \item \textbf{Выручка}: \texttt{Sum(F('unit\_price') * F('quantity'))};
    \item \textbf{Себестоимость}: \texttt{Sum(F('component\_\allowbreak purchase\_price') * F('quantity'))};
    \item \textbf{Прибыль}: разность выручки и себестоимости;
    \item \textbf{Количество заказов} и \textbf{средний чек}: \texttt{Count('order', distinct=True)} и отношение выручки к количеству заказов.
\end{itemize}

\textbf{Топ-10 товаров} рассчитывается группировкой позиций заказов по \texttt{component\_\allowbreak title} и \texttt{component\_\allowbreak brand} с агрегацией \texttt{Sum('quantity')} и сортировкой по убыванию. Дополнительно выгружаются позиции с критически низкими остатками и подсчитывается количество товаров «нет в наличии».

\textbf{Визуализация} построена на цветных KPI-карточках (с градиентным фоном для основных показателей), таблице топ-товаров и боковом списке низких остатков с цветными бейджами (жёлтый --- мало, красный --- нет в наличии).

\subsection{Реализация загрузки и выгрузки файлов}

Модуль импорта-экспорта реализован для каталога комплектующих и отчётности по заказам.

\textbf{Экспорт каталога} (\texttt{ComponentExportView}) --- генерирует CSV-файл с разделителем «;», кодировкой UTF-8 и предваряющим BOM (для корректного открытия в Microsoft Excel). Заголовок: \texttt{id, category, title, brand, model, specs, purchase\_price, sale\_price, stock, is\_active}. Содержимое поля \texttt{specs} (JSON) сериализуется в текстовый JSON через \texttt{json.dumps(ensure\_ascii=False)}.

\textbf{Импорт каталога} (\texttt{ComponentImportView}) --- загрузка CSV-файла через стандартный \texttt{FileField}. Алгоритм импорта:
\begin{enumerate}
    \item декодирование содержимого как \texttt{utf-8-sig} (корректная обработка BOM из Excel);
    \item чтение через \texttt{csv.DictReader} c разделителем «;»;
    \item для каждой строки --- валидация обязательных полей (категория, название);
    \item поиск или автоматическое создание категории по \texttt{slug};
    \item парсинг JSON-поля \texttt{specs};
    \item операция \texttt{Component.objects.update\_or\_create(title=...)}, что обеспечивает идемпотентность повторных загрузок.
\end{enumerate}

По завершении импорта пользователю выводится сообщение со статистикой (создано, обновлено, пропущено) и до 10 ошибок валидации в виде flash-сообщений.

\textbf{Экспорт заказов} (\texttt{OrderExportView}) --- CSV-выгрузка заказов за период с заголовком \texttt{order\_id, date, client\_email, status, total\_sum, payment\_method, delivery\_method}. Период фильтруется параметрами \texttt{?from=YYYY-MM-DD\&to=YYYY-MM-DD}.

\subsection*{Выводы по главе 3}
\addcontentsline{toc}{subsection}{Выводы по главе 3}

В главе описана программная реализация всех шести задач раздела~3 пояснительной записки. Реализованы базовая структура приложения и шаблоны страниц на Bootstrap~5 с современным дизайн-кодом, аутентификация по email и ролевая модель доступа, полноценные CRUD-интерфейсы для всех основных сущностей системы, интерактивный конфигуратор сборки ПК с проверкой совместимости в реальном времени и расчётом потребляемой мощности, аналитический модуль с расчётом выручки, прибыли, среднего чека и топ-товаров, а также модуль импорта и экспорта файлов в формате CSV. Полученная реализация полностью покрывает функциональные требования, сформулированные в подразделе~1.4.

Исходный код проекта размещён в открытом репозитории:
\begin{center}
    \url{https://github.com/Haskeri/kpweb}
\end{center}

Развёрнутое приложение доступно по адресу:
\begin{center}
    \url{https://exam.anufriev.pvpcult.ru/}
\end{center}
